
    %%%%%%%%%%%%%%%%%%%%%%%%%%%%%%%%%%%%%%%%%%%%%%%%%%%%%%%%%%%%%%%%%%%%%%%%%%%%%%%%
    %%% Classe do documento
    \documentclass[12pt, addpoints]{exam}
    \UseRawInputEncoding

    %%%%%%%%%%%%%%%%%%%%%%%%%%%%%%%%%%%%%%%%%%%%%%%%%%%%%%%%%%%%%%%%%%%%%%%%%%%%%%%%
    % preambulo.tex
    %
    % Arquivo de configuração de packages para uso o LaTeX, conforme minhas
    % preferências e modelos pessoais.
    %
    % Para maiores informações, visite:
    %    https://github.com/abrantesasf/latex
    %
    % NÃO ALTERE SE NÃO SOUBER O QUE ESTÁ FAZENDO!
    %%%%%%%%%%%%%%%%%%%%%%%%%%%%%%%%%%%%%%%%%%%%%%%%%%%%%%%%%%%%%%%%%%%%%%%%%%%%%%%%


    %%%%%%%%%%%%%%%%%%%%%%%%%%%%%%%%%%%%%%%%%%%%%%%%%%%%%%%%%%%%%%%%%%%%%%%%%%%%%%%%
    %%% Estruturas de controle:
    \input{static/utils/controles}

    %%%%%%%%%%%%%%%%%%%%%%%%%%%%%%%%%%%%%%%%%%%%%%%%%%%%%%%%%%%%%%%%%%%%%%%%%%%%%%%%
    %%% Compilação condicional em PDF:
    \input{static/utils/pdf}

    %%%%%%%%%%%%%%%%%%%%%%%%%%%%%%%%%%%%%%%%%%%%%%%%%%%%%%%%%%%%%%%%%%%%%%%%%%%%%%%%
    %%% Configurações de layout da página:
    \input{static/utils/layout}

    %%%%%%%%%%%%%%%%%%%%%%%%%%%%%%%%%%%%%%%%%%%%%%%%%%%%%%%%%%%%%%%%%%%%%%%%%%%%%%%%
    %%% Configurações para autores:
    \input{static/utils/autores}

    %%%%%%%%%%%%%%%%%%%%%%%%%%%%%%%%%%%%%%%%%%%%%%%%%%%%%%%%%%%%%%%%%%%%%%%%%%%%%%%%
    %%% Cabeçalho e rodapé:
    %%% Atenção! É MUITO DIFÍCIL ajustar apropriadamente os running headers
    %%% e running footers da primeira vez. Você pode confiar no LaTeX e deixar que
    %%% ele faça o ajuste automaticamente ou pode quebrar a cabeça e
    %%% tentar melhorar algo que o LaTeX já faz muito bem, mesmo que não fique
    %%% exatamente do jeito que você quer. O que você prefere? Se quiser ajustar
    %%% manualmente, descomente a linha abaixo e faça as configurações no arquivo
    %%% "utils/headings.tex".
    %\input{static/utils/headings}

    %%%%%%%%%%%%%%%%%%%%%%%%%%%%%%%%%%%%%%%%%%%%%%%%%%%%%%%%%%%%%%%%%%%%%%%%%%%%%%%%
    %%% Configuração para as fontes do documento:
    %%% Se você quiser usar fontes próprias ou do sistema, precisará ajustar as
    %%% configurações no arquivo "utils/fontes.tex".
    %%% ATENÇÃO: por padrão eu utilizo XeLaTeX com fontes proprietárias projetadas
    %%% por Matthew Butterick, adquiridas em https://mbtype.com, que não podem ser
    %%% distribuídas devido à licença de utilização. Se você usar XeLaTeX, você
    %%% DEVERÁ OBRIGATORIAMENTE ajustar as fontes no arquivo "utils/fontes.tex".
    \input{static/utils/fontes}

    %%%%%%%%%%%%%%%%%%%%%%%%%%%%%%%%%%%%%%%%%%%%%%%%%%%%%%%%%%%%%%%%%%%%%%%%%%%%%%%%
    %%% Sumário:
    \input{static/utils/toc}

    %%%%%%%%%%%%%%%%%%%%%%%%%%%%%%%%%%%%%%%%%%%%%%%%%%%%%%%%%%%%%%%%%%%%%%%%%%%%%%%%
    %%% Matemática:
    \input{static/utils/matematica}

    %%%%%%%%%%%%%%%%%%%%%%%%%%%%%%%%%%%%%%%%%%%%%%%%%%%%%%%%%%%%%%%%%%%%%%%%%%%%%%%%
    %%% Notas de rodapé:
    \input{static/utils/footnotes}

    %%%%%%%%%%%%%%%%%%%%%%%%%%%%%%%%%%%%%%%%%%%%%%%%%%%%%%%%%%%%%%%%%%%%%%%%%%%%%%%%
    %%% Referências cruzadas, links, citações:
    \input{static/utils/crossrefs}

    %%%%%%%%%%%%%%%%%%%%%%%%%%%%%%%%%%%%%%%%%%%%%%%%%%%%%%%%%%%%%%%%%%%%%%%%%%%%%%%%
    %%% Configurações para os hiperlinks em PDF:
    \input{static/utils/pdfconfig}

    %%%%%%%%%%%%%%%%%%%%%%%%%%%%%%%%%%%%%%%%%%%%%%%%%%%%%%%%%%%%%%%%%%%%%%%%%%%%%%%%
    %%% Glossário e índice remissivo:
    \input{static/utils/glossindex}

    %%%%%%%%%%%%%%%%%%%%%%%%%%%%%%%%%%%%%%%%%%%%%%%%%%%%%%%%%%%%%%%%%%%%%%%%%%%%%%%%
    %%% Referências bibliográficas:
    \input{static/utils/bibliografia}

    %%%%%%%%%%%%%%%%%%%%%%%%%%%%%%%%%%%%%%%%%%%%%%%%%%%%%%%%%%%%%%%%%%%%%%%%%%%%%%%%
    %%% Cores:
    \input{static/utils/cores}

    %%%%%%%%%%%%%%%%%%%%%%%%%%%%%%%%%%%%%%%%%%%%%%%%%%%%%%%%%%%%%%%%%%%%%%%%%%%%%%%%
    %%% Computação:
    \input{static/utils/computacao}

    %%%%%%%%%%%%%%%%%%%%%%%%%%%%%%%%%%%%%%%%%%%%%%%%%%%%%%%%%%%%%%%%%%%%%%%%%%%%%%%%
    %%% Gráficos:
    \input{static/utils/graficos}

    %%%%%%%%%%%%%%%%%%%%%%%%%%%%%%%%%%%%%%%%%%%%%%%%%%%%%%%%%%%%%%%%%%%%%%%%%%%%%%%%
    %%% Ícones:
    \input{static/utils/icones}

    %%%%%%%%%%%%%%%%%%%%%%%%%%%%%%%%%%%%%%%%%%%%%%%%%%%%%%%%%%%%%%%%%%%%%%%%%%%%%%%%
    %%% Blocos:
    \input{static/utils/blocos}

    %%%%%%%%%%%%%%%%%%%%%%%%%%%%%%%%%%%%%%%%%%%%%%%%%%%%%%%%%%%%%%%%%%%%%%%%%%%%%%%%
    %%% Boxes coloridos:
    \input{static/utils/colorbox}

    %%%%%%%%%%%%%%%%%%%%%%%%%%%%%%%%%%%%%%%%%%%%%%%%%%%%%%%%%%%%%%%%%%%%%%%%%%%%%%%%
    %%% Tabelas:
    \input{static/utils/tabelas}

    %%%%%%%%%%%%%%%%%%%%%%%%%%%%%%%%%%%%%%%%%%%%%%%%%%%%%%%%%%%%%%%%%%%%%%%%%%%%%%%%
    %%% Ambientes de listas:
    \input{static/utils/listas}

    %%%%%%%%%%%%%%%%%%%%%%%%%%%%%%%%%%%%%%%%%%%%%%%%%%%%%%%%%%%%%%%%%%%%%%%%%%%%%%%%
    %%% Ambientes floats e similares:
    \input{static/utils/floats}

    %%%%%%%%%%%%%%%%%%%%%%%%%%%%%%%%%%%%%%%%%%%%%%%%%%%%%%%%%%%%%%%%%%%%%%%%%%%%%%%%
    %%% Ajustes para a classe EXAM:
    \input{static/utils/exam}

    %%%%%%%%%%%%%%%%%%%%%%%%%%%%%%%%%%%%%%%%%%%%%%%%%%%%%%%%%%%%%%%%%%%%%%%%%%%%%%%%
    %%% Ajustes para textos:
    \input{static/utils/texto}

    %%%%%%%%%%%%%%%%%%%%%%%%%%%%%%%%%%%%%%%%%%%%%%%%%%%%%%%%%%%%%%%%%%%%%%%%%%%%%%%%
    %%% Ambientes, aliases, comandos e customizações em geral para serem
    %%% utilizadas nos documentos. Defina aqui qualquer customização específica
    %%% para seus documentos!
    \input{static/utils/customizacoes}

    %%%%%%%%%%%%%%%%%%%%%%%%%%%%%%%%%%%%%%%%%%%%%%%%%%%%%%%%%%%%%%%%%%%%%%%%%%%%%%%%
    %%% Ajuste do layout, espaçamento de linhas e etc.:
    \geometry{a4paper, portrait, left=2.0cm, right=2.0cm, top=2.0cm, bottom=2.0cm}
    \singlespacing

    %%%%%%%%%%%%%%%%%%%%%%%%%%%%%%%%%%%%%%%%%%%%%%%%%%%%%%%%%%%%%%%%%%%%%%%%%%%%%%%%
    %%% Configurações de cabeçalho e rodapé (não use fancyhdr pois ocorre conflito):
    \pagestyle{headandfoot}
    \runningheadrule
    \firstpageheader{}{}{}
    \runningheader{Sem disciplina}{}{Avaliação}
    \firstpagefooter{}{}{Página \thepage\ de \numpages}
    \runningfooter{}{}{Página \thepage\ de \numpages}
    \runningfootrule

    %%%%%%%%%%%%%%%%%%%%%%%%%%%%%%%%%%%%%%%%%%%%%%%%%%%%%%%%%%%%%%%%%%%%%%%%%%%%%%%%
    %%% Configurações para as propriedades do PDF:
    \hypersetup{
    hidelinks,           % Comente para web, descomente para impressão
    colorlinks=true,      % True para web, False para impressão
    pdftitle=GAGA: Avaliação,
    pdfauthor=Matheus Endlich Silveira,
    pdfsubject=Sem disciplina,
    pdfkeywords=['Sem tag'],
    pdfinfo={
        CreationDate={}, % Ex.: D:AAAAMMDDHH24MISS
        ModDate={}       % Ex.: D:AAAAMMDDHH24MISS
    }
    }
    \input{static/utils/pdfconfig}

    %%%%%%%%%%%%%%%%%%%%%%%%%%%%%%%%%%%%%%%%%%%%%%%%%%%%%%%%%%%%%%%%%%%%%%%%%%%%%%%%
    %%% Começa o documento
    \begin{document}

    %%%%%%%%%%%%%%%%%%%%%%%%%%%%%%%%%%%%%%%%%%%%%%%%%%%%%%%%%%%%%%%%%%%%%%%%%%%%%%%%
    %%% Folha de instruções padronizadas da UVV para a prova
    \pdfbookmark[1]{Instruções}{instrucoes}
    \begin{figure}[H]
        \begin{center}
            \includegraphics[scale=0.3]{{static/imgs/logo-uvv.png}}
        \end{center}	
    \end{figure}

    \vspace{-1cm}

    \begin{table}[H]
        \centering
        \begin{tabular}{|llll|}
            \hline
            \multicolumn{2}{|l|}{\textbf{Disciplina:} Sem disciplina} & 
            \multicolumn{1}{l|}{\multirow{3}{*}{\textbf{\makecell{Nota:\\ \,\\ \,}}}} & 
            \multirow{3}{*}{\textbf{\makecell{Coordenador:\\ \,\\ \,}}} \\ 
            \cline{1-2}
            \multicolumn{2}{|l|}{\textbf{Professor:} Matheus Endlich Silveira \hspace{4.72cm} } & 
            \multicolumn{1}{l|}{} & \\ 
            \cline{1-2}
            \multicolumn{2}{|l|}{\textbf{Aluno:}} & 
            \multicolumn{1}{l|}{} & \\ 
            \hline
            \multicolumn{1}{|l|}{\textbf{Turma:} \hspace{6.3cm} } & 
            \multicolumn{1}{l|}{\textbf{Semestre:}} & 
            \multicolumn{2}{l|}{\textbf{Valor:} 10 pontos} \\ 
            \hline
            \multicolumn{1}{|l|}{\textbf{Data:}} & 
            \multicolumn{3}{l|}{\textbf{Avaliação:}} \\ 
            \hline
        \end{tabular}
    \end{table}

    \begin{center}
        \fbox{\setlength\fboxsep{0.3cm} \fbox{\parbox{15.4cm}{
            \begin{center}
                \textbf{INSTRUÇÕES PARA A REALIZAÇÃO DESTA AVALIAÇÃO:}
            \end{center}
            \begin{itemize}[leftmargin=*]
                \item Esta é a primeira avaliação da disciplina Sem disciplina, 
                    e engloba o conteúdo referente Sem tag.
                \item Esta avaliação contém 50 questões objetivas, \textbf{todas obrigatórias}.
                \item \textbf{Leia atentamente cada questão!} As questões objetivas terão uma,
                    e apenas uma única, resposta a ser marcada.
                \item \textbf{Desligue o celular} antes de começar. A avaliação é
                    \textbf{individual} e \textbf{sem consulta}.
                \item As questões podem ser respondidas, na prova, com lápis ou caneta. Ao final
                    da prova \textbf{{VOCÊ DEVERÁ PREENCHER O GABARITO COM CANETA
                    DE COR PRETA}} \textbf{\underline{OU AZUL}} para que seja feita a correção
                    automatizada através da leitura do padrão de respostas marcado no
                    gabarito.
                \item \textbf{CUIDADO ao preencher o gabarito} pois eles são \textbf{INDIVIDUAIS
                    e NOMEADOS} (cada estudante tem um gabarito próprio específico, com um
                    código de identificação único). \textbf{Questões rasuradas são anuladas}
                    pelo software de leitura do gabarito.
                \item Ao final da prova, \textbf{DEVOLVA A PROVA E O GABARITO} para o professor.
                \item Siga todas as normas de \textbf{Integridade Acadêmica} da disciplina.
                    Alunos que forem flagrados com qualquer espécie de ``cola'' ou trocando
                    informações com outros alunos terão suas avaliações recolhidas, as notas
                    zeradas, e a situação será encaminhada para a coordenação para a aplicação
                    das penalidades previstas pela universidade.
                \item Com 90 minutos de avaliação você terá tempo de até 2,min por questão,
                    em média.
                \item Boa avaliação!
            \end{itemize}
            \vspace{2cm}
        }}}
    \end{center}
    \newpage

    %%%%%%%%%%%%%%%%%%%%%%%%%%%%%%%%%%%%%%%%%%%%%%%%%%%%%%%%%%%%%%%%%%%%%%%%%%%%%%%%
    %%% Inicia o bloco de questões:
    \begin{questions}

    \newpage
    \question

O contador da figura abaixo é:

\begin{figure}

    \centering

    \includegraphics[width=0.5\linewidth]{static/imgs/defaul.png}

    \label{fig:enter-label}

\end{figure}
\begin{choices}
\choice isócrono
\choice auto-sincronizado
\choice assíncrono
\choice anisócrono
\choice síncrono
\end{choices}

\question Assinale a alternativa \textbf{incorreta}:
\begin{choices}
\choice O MariaDB é um SGBD
\choice O modelo relacional de dados foi criado por Edgar Frank Cood
\choice Um SGBD não tem interface gráfica
\choice Modelo de dados é a mesma coisa que o diagrama relacional
\choice No modelo relacional, tudo o que o usuário percebe são tabelas
\end{choices}

\question

Considere as seguintes afirmações:

\begin{enumerate}[label=\Roman*.]

    \item A classe de linguagens aceita por um Autômato Finito Determinístico (AFD) não é a mesma que um Autômato Finito Não Determinístico (AFND).

    \item Para algumas expressões regulares não é possível construir um AFD.

    \item A expressão regular \((b + ba)^+\) aceita as "strings" de b's e a's começando com b e não tendo dois a's consecutivos.

\end{enumerate}

Selecione a afirmativa correta:
\begin{choices}
\choice As afirmativas I e III são falsas.
\choice Apenas a afirmativa III é verdadeira.
\choice As afirmativas I e III são verdadeiras.
\choice As afirmativas I e II são verdadeiras.
\choice As afirmativas II e III são falsas.
\end{choices}

\question

Quando dizemos que um banco de dados é em \ingles{cloud} e centralizado, estamos

nos referindo, respectivamente, à localização dos dados e da

infraestrutura. Verdadeiro ou falso?
\begin{choices}
\choice Não é possível determinar com os dados apresentados.
\choice Verdadeiro
\choice Falso
\end{choices}

\question

Seu chefe que fazer uma campanha de venda para produtos que custam entre 20 e 30

reais, mas está preocupado porque ouviu de um dos vendedores que diversos

produtos nessa faixa de preço estão com o estoque zerado. Para verificar a

situação seu chefe pediu para que você prepare um relatório que mostre os

produtos que custam entre 20 e 30 reais e que estão com o estoque zerado. Você

preparou o seguinte relatório para seu chefe:

\begin{tcolorbox}

 \begin{verbatim}1  SELECT l.nome AS loja

2       , p.nome AS produto

3       , p.preco_unitario AS preco

4       , e.quantidade AS qtd

5  FROM produtos p

6  INNER JOIN estoques e ON (e.produto_id = p.produto_id)

7  INNER JOIN lojas l ON (e.loja_id = l.loja_id)

8  WHERE (p.preco_unitario >= 20 OR p.preco_unitario <= 30)

9    AND e.quantidade = 0

10 ORDER BY loja, produto;

 \end{verbatim}

\end{tcolorbox}

O seu relatório:
\begin{choices}
\choice Está errado pois a tabela que deveria estar na cláusula FROM é a tabela ``lojas', não a tabela ``produtos'.
\choice Está errado pois vai trazer também produtos fora da faixa de preço que o seu chefe quer.
\choice Nenhuma das alternativas anteriores.
\choice Está errado pois há um erro de sintaxe nas linhas 6 e 7.
\choice Está correto e trará exatamente, os produtos com o estoque zerado e na faixa de preço requisitada pelo seu chefe. Além disso ordena o resultado pelo nome da loja e pelo nome do produto.
\end{choices}

\question

A velocidade de um ponto em movimento é dada pela equação\\

$v(t) = t e^{-0.01t} \, \text{m/s}$\\

O espaço percorrido desde o instante que o ponto começou a se mover até a sua parada total é


\begin{choices}
\choice \( 10^2 \, m \)
\choice \( (e^{-100} - 1) m \)
\choice \( 10^4 \, m \)
\choice \( 10^3 e^{-0.01} \, m \)
\choice \( 10^2 e^{-1} \, m \)
\end{choices}

\question

(ENADE 2014) Considere o seguinte banco de dados ilustrado no diagrama relaconal

abaixo:

\begin{figure}[H]

  \centering

  \caption{Estados e fornecedores}

  \label{fig:estados}

  %\fbox{

    \includegraphics[width=0.5\linewidth]{static/imgs/defaul.png}

  %}\\

  %\footnotesize{Fonte: xxxx.}

\end{figure}

A expressão SQL que obtém os nomes dos estados para os quais não há fornecedores

cadastrados é:


\begin{choices}
\choice \begin{verbatim}SELECT e.nome

FROM estado e,

FROM fornecedor f

WHERE e.uf = f.uf;

\end{verbatim}


\choice \begin{verbatim}SELECT e.nome

FROM estado e

WHERE e.uf IN (SELECT f.uf

               FROM fornecedor f);

\end{verbatim}


\choice \begin{verbatim}SELECT e.nome

FROM estado e,

FROM fornecedor f

WHERE e.nome = f.uf;

\end{verbatim}


\choice \begin{verbatim}SELECT e.uf

FROM estado e

WHERE e.nome NOT IN (SELECT f.uf

                     FROM fornecedor f);

\end{verbatim}
\choice \begin{verbatim}SELECT e.nome

FROM estado e

WHERE e.uf NOT IN (SELECT f.uf

                   FROM fornecedor f);

\end{verbatim}


\end{choices}

\question

Considere um grafo \( G \) satisfazendo as seguintes propriedades:



\begin{itemize}

    \item \( G \) é conexo

    \item Se removermos qualquer aresta de \( G \), o grafo obtido é desconexo.

\end{itemize}



Então é correto afirmar que o grafo \( G \) é:
\begin{choices}
\choice Hamiltoniano
\choice Uma árvore
\choice Euleriano
\choice Não bipartido
\choice Um circuito
\end{choices}

\question

\textbf Supondo a Relação \texttt{PROJ (PNO, Nome, Orçam)}, com chave primária \texttt{PNO} e a Relação \texttt{DSG (ENO, PNO, Dur, Resp)}, com chave primária \{\texttt{ENO, PNO}\} e chave estrangeira \texttt{PNO} em relação a \texttt{PROJ}, a asserção abaixo \textbf{NÃO} expressa:



\begin{equation}\forall g \in \texttt{DSG}, \exists j \in \texttt{PROJ} : g.\texttt{PNO} = j.\texttt{PNO}\end{equation}
\begin{choices}
\choice Uma restrição de integridade de chave primária em $\texttt{PROJ}$.
\choice Uma restrição a ser verificada na inserção de tuplas em $\texttt{DSG}$.
\choice Uma restrição que define um estado consistente do banco de dados.
\choice Uma restrição a ser verificada na atualização de tuplas em $\texttt{DSG}$.
\choice Uma restrição de integridade de chave estrangeira em $\texttt{DSG}$.
\end{choices}

\question

Uma árvore binária é declarada em C como



\begin{verbatim}

typedef struct no *apontador;

struct no {

    int valor;

    apontador esq, dir;

};

\end{verbatim}



onde \texttt{esq} e \texttt{dir} representam ligações para os filhos esquerdo e direito de um nó da árvore, respectivamente. Qual das seguintes alternativas é uma implementação correta da operação que inverte as posições dos filhos esquerdo e direito de um nó \texttt{p} da árvore, onde \texttt{t} é um apontador auxiliar.


\begin{choices}
\choice \texttt{p->esq = p->dir;} \\

          \texttt{t = p->esq;} \\

          \texttt{p->dir = t;}
\choice \texttt{p->dir = t;} \\

          \texttt{p->esq = p->dir;} \\

          \texttt{p->dir = t;}
\choice \texttt{t = p;} \\

          \texttt{p->esq = p->dir;} \\

          \texttt{p->dir = p->esq;}
\choice \texttt{t = p->dir;} \\

          \texttt{p->esq = p->dir;} \\

          \texttt{p->dir = t;}
\choice \texttt{t = p->dir;} \\

          \texttt{p->dir = p->esq;} \\

          \texttt{p->esq = t;}
\end{choices}

\question

Em relação aos metadados de um banco de dados, assinale a afirmação correta:
\begin{choices}
\choice São atualizados pelo SGBD, sempre que ocorre alguma operação de inserção, atualização ou remoção de dados.
\choice Como são sempre atualizados de forma automática pelo sistema de gerenciamento de bancos de dados, estão sempre prontos para nos dar uma ``foto' da estrutura do banco de dados em um momento no tempo.
\choice São armazenados no catálogo (ou dicionário de dados) do sistema de gerenciamento de bancos de dados, que fica nas mesmas tabelas de dados dos usuários, permitindo assim que os metadados estejam sempre atualizados.
\choice Armazenam informações como os dados de cadastro dos usuários, os valores de compra em uma ordem de compra, as turmas e disciplinas que os professores dão aula, e dados de recursos humanos.
\choice Armazenam informações sobre os dados dos usuários, ou seja, são dados de rotina dos usuários, aqueles dados que os usuários precisam para trabalhar no seu dia a dia.
\end{choices}

\question

\begin{figure}[H]

    \centering

    \includegraphics[width=0.5\linewidth]{static/imgs/defaul.png}

    \caption{Enter Caption}

\end{figure}

Se \( O = (0, 0, 0) \); \( A = (2, 4, 1) \); \( B = (3, 1, 1) \) e \( C = (1, 3, 5) \) então o volume do sólido acima é


\begin{choices}
\choice 30
\choice 21
\choice 44
\choice 35
\choice 35/2
\end{choices}

\question

Dentre as definições a seguir, ligadas ao conceito de normalização do modelo relacional, qual delas é \textbf{INCORRETA}?


\begin{choices}
\choice As formas normais se baseiam em certas estruturas de dependências.
\choice A primeira forma normal estabelece que os atributos da relação contêm apenas valores atômicos.
\choice A normalização é um processo passo a passo reversível de substituição de uma dada coleção de relações por sucessivas coleções de relações as quais possuem uma estrutura progressivamente mais simples e mais regular.
\choice O objetivo da normalização é eliminar várias anomalias (ou aspectos indesejáveis) de uma relação.
\choice As relações que obedecem à primeira forma normal não apresentam anomalias.
\end{choices}

\question

Considere a seguinte tabela em uma base de dados relacional (chave primária sublinhada):\\

Tabela1 (CodAluno, CodDisciplina, AnoSemestre, NomeAluno, NomeDisciplina, CodNota, DescricaoNota)\\

Considere as seguintes dependências funcionais:

\begin{itemize}

    \item CodAluno \(\rightarrow\) NomeAluno

    \item CodDisciplina \(\rightarrow\) NomeDisciplina

    \item (CodAluno, CodDisciplina, AnoSemestre) \(\rightarrow\) CodNota

    \item (CodAluno, CodDisciplina, AnoSemestre) \(\rightarrow\) DescricaoNota

    \item CodNota \(\rightarrow\) DescricaoNota

\end{itemize}



Considerando as formas normais, qual das afirmativas abaixo se aplica:
\begin{choices}
\choice A tabela encontra-se na primeira forma normal, mas não na segunda forma normal.
\choice A tabela encontra-se na segunda forma normal, mas não na terceira forma normal.
\choice A tabela encontra-se na terceira forma normal, mas não na quarta forma normal.
\choice A tabela está na quarta forma normal.
\choice A tabela não está na primeira forma normal.
\end{choices}

\question

O Modelo Espiral de desenvolvimento de software, proposto por Barry Boehm na década de 1970, apresenta as seguintes afirmações:

\begin{enumerate}[label=\Roman*.]

    \item Suas atividades de desenvolvimento são conduzidas por riscos;

    \item Cada ciclo da espiral inclui 4 passos: passo 1 - identificação dos objetivos; passo 2 - avaliação das alternativas tendo em vista os objetivos e os riscos (incertezas, restrições) do desenvolvimento; passo 3 - desenvolvimento de estratégias (simulação, prototipagem) para resolver riscos; e passo 4 - planejamento do próximo passo e continuidade do processo determinada pelos riscos restantes;

    \item É um modelo evolutivo em que cada passo pode ser representado por um quadrante num diagrama cartesiano: assim, na dimensão radial da espiral tem-se o custo acumulado dos vários passos do desenvolvimento, enquanto na dimensão angular tem-se o progresso do projeto.

\end{enumerate}

Levando-se em conta as três afirmações I, II e III acima, identifique a única alternativa válida:


\begin{choices}
\choice Apenas a I e a II estão corretas;
\choice Apenas a II e a III estão corretas;
\choice Apenas a III está correta.
\choice Apenas a I e a III estão corretas;
\choice As afirmações I, II e III estão corretas;
\end{choices}

\question

Considere as seguintes afirmações sobre redes neurais artificiais: \begin{enumerate}[label=\Roman*.] \item Um perceptron elementar só computa funções linearmente separáveis. \item Não aceitam valores numéricos como entrada. \item O "conhecimento" é representado principalmente através do peso das conexões. \end{enumerate} São corretas:
\begin{choices}
\choice I, II e III
\choice Apenas III
\choice Apenas II e III
\choice Apenas I e III
\choice Apenas I e II 
\end{choices}

\question

Por que não podemos utilizar apelidos de coluna na cláusula WHERE da SQL?
\begin{choices}
\choice Porque só podemos utilizar aliases, e não apelidos.
\choice Nós podemos utilizar apelidos na cláusula WHERE.
\choice Porque a cláusula WHERE é avaliada antes da SELECT.
\choice Porque os apelidos só são avaliados depois da cláusula ORDER BY.
\choice Porque a cláusula SELECT é avaliada antes da WHERE.
\end{choices}

\question

A média aritmética de uma lista de 50 números é 50. Se dois desses números, 51 e 97, forem suprimidos dessa lista a média dos restantes será


\begin{choices}
\choice 49
\choice 50
\choice 51
\choice 47
\choice 40
\end{choices}

\question

Considere as seguintes tabelas em uma base de dados relacional (chaves primárias sublinhadas):\\

Departamento (\underline{CodDepto}, NomeDepto)\\

Empregado (\underline{CodEmp}, NomeEmp, CodDepto)\\

Considere as seguintes restrições de integridade sobre esta base de dados relacional:\\

- Empregado.CodDepto é sempre diferente de NULL\\

- Empregado.CodDepto é chave estrangeira da tabela Departamento com cláusulas ON DELETE RESTRICT e ON UPDATE RESTRICT\\

Qual das seguintes validações não é especificada por estas restrições de integridade:
\begin{choices}
\choice Sempre que o valor de Departamento.CodDepto for alterado, deve ser garantido que não há uma linha com o antigo valor de Departamento.CodDepto na coluna Empregado.CodDepto.
\choice Sempre que uma linha for excluída de Departamento, deve ser garantido que o valor de Departamento.CodDepto não aparece na coluna Empregado.CodDepto.
\choice Sempre que uma nova linha for inserida em Empregado, deve ser garantido que o valor de Empregado.CodDepto aparece na coluna Departamento.CodDepto.
\choice Sempre que uma nova linha for inserida em Departamento, deve ser garantido que o valor de Departamento.CodDepto aparece na coluna Empregado.CodDepto.
\choice Sempre que o valor de Empregado.CodDepto for alterado, deve ser garantido que o novo valor de Empregado.CodDepto aparece em Departamento.CodDepto.
\end{choices}

\question

Para que serve a segmentação de um processador (\textit{pipelining})?
\begin{choices}
\choice aumentar a velocidade do relógio
\choice simplificar o conjunto de instruções
\choice simplificar a implementação do processador
\choice reduzir o número de instruções estáticas nos programas
\choice permitir a execução de mais de uma instrução por ciclo de relógio
\end{choices}

\question

Pode-se afirmar que o gráfico da função \(y = 2 + \frac{1}{x - 1}\) é o gráfico da função \(y = \frac{1}{x}\):
\begin{choices}
\choice transladado uma unidade para a direita e duas unidades para cima
\choice transladado uma unidade para a direita e duas unidades para baixo
\choice transladado uma unidade para a esquerda e duas unidades para cima
\choice nenhuma das anteriores.
\choice transladado uma unidade para a esquerda e duas unidades para baixo
\end{choices}

\question

Qual das seguintes opções \textbf{não é} uma vantagem do uso de um banco de

dados sobre o sistema de arquivos convencional?


\begin{choices}
\choice Estruturação de dados
\choice Controle de redundância e inconsistência
\choice Eficiência
\choice Controle de acesso e visualização de dados
\choice Independência de dados
\end{choices}

\question

Sistemas de processamento de transações, tais como sistemas de reservas aéreas, devem prover um mecanismo que garanta que cada transação não é afetada por outras transações que possam estar ocorrendo ao mesmo tempo. Transações de duas fases obedecem a um protocolo que garante essa atomicidade. Em transações de duas fases:


\begin{choices}
\choice Verifica-se a disponibilidade de todas as travas antes de executar qualquer ação de travamento.
\choice Todas as operações de leitura ocorrem antes da primeira operação de escrita.
\choice Uma trava compartilhada sobre um objeto deve ser obtida antes de uma trava exclusiva sobre o objeto ser obtida.
\choice Qualquer objeto corretamente travado deve ser destravado antes que outro objeto possa ser travado.
\choice Todas as ações de travamento (\textit{lock}) ocorrem antes da primeira ação de destravamento.
\end{choices}

\question

Em qual das situações abaixo um sistema de Raciocínio Baseado em Casos não deve ser utilizado?


\begin{choices}
\choice Quando for fácil a obtenção de regras
\choice Quando a experiência for tão valiosa quanto o conhecimento em livros texto.
\choice Quando as regras utilizadas apresentam um grande número de exceções.
\choice Em aplicações de diagnóstico médico.
\choice Quando especialistas conversam sobre seus domínios dando exemplos.
\end{choices}

\question

Considere a seguinte tabela para uma base de dados relacional:\\

Empregado (CodEmp, NomeEmp, CodDepto)\\

Considere que esta tabela tem um índice na forma de uma árvore B sobre as colunas (CodEmp, CodDepto), nesta ordem.

Quanto a este índice, considere as seguintes afirmativas:

\begin{enumerate}

    \item Este índice pode ser usado pelo SGBD relacional para acelerar uma consulta na qual são fornecidos os valores de CodEmp e CodDepto.

    \item Este índice pode ser usado pelo SGBD relacional para acelerar uma consulta na qual é fornecido um valor de CodEmp.

    \item Este índice não é adequado para ser usado pelo SGBD relacional para acelerar uma consulta na qual é fornecido um valor de CodDepto.

    \item O algoritmo que faz inserções e remoções de entradas do índice tem por objetivo garantir que o índice fique organizado de tal forma que o acesso a cada nodo da árvore implique em número de acessos semelhantes.

    \item O índice por árvore-B não é adequado para tabelas que sofrem grande número de inclusões e exclusões, pois exige reorganizações frequentes.

\end{enumerate}





Quanto a estas afirmativas pode se dizer que:
\begin{choices}
\choice Nenhuma das afirmativas está correta
\choice Apenas as afirmativas 1), 2) e 4) estão corretas
\choice Apenas as afirmativas 1), 2), 3) e 4) estão corretas
\choice Todas afirmativas estão corretas
\choice Apenas as afirmativas 1), 2) e 5) estão corretas
\end{choices}

\question

Considere as seguintes afirmações sobre mecanismos de inferência em sistemas baseados em regras.

\begin{enumerate}[label=\Roman*.]

    \item O encadeamento regressivo tem pouca utilidade prática, pois deve partir do possível resultado.

    \item O encadeamento progressivo tanto pode ser em amplitude quanto em profundidade.

    \item Podem trabalhar com informações incertas ou incompletas. 

\end{enumerate}

São corretas:
\begin{choices}
\choice Apenas III
\choice Apenas II e III
\choice I, II e III
\choice Apenas I e III
\choice Apenas I e II
\end{choices}

\question

No programa abaixo, escrito em Pascal, os parâmetros do procedimento \texttt{vr} são passados por valor.



\begin{verbatim}

program teste;

var x, y: integer;



procedure vr(u, v: integer);

begin

    u := 2 * u;

    x := u + v;

    u := u - 1;

end;



begin

    x := 4;

    y := 2;

    vr(x, y);

    writeln(x);

end.

\end{verbatim}



O valor de \texttt{x} impresso na última linha do programa é:
\begin{choices}
\choice 10
\choice 8
\choice 5
\choice 4
\choice 7
\end{choices}

\question

O usuário que não é da área de tecnologia, costuma ser de alta gerência, que é

especialista no negócio e tem a responsabilidade central pelas políticas de

dados da empresa é o:
\begin{choices}
\choice Gerente de Tecnologia da Informação (ITM: IT Manager)
\choice Administrador de Dados (DA: Data Administrator)
\choice Administrador do Banco de Dados (DBA: Database Administrator)
\choice Coordenador de Tecnologia da Informação (ITS: IT Supervisor)
\choice Encarregado da Proteção de Dados (DPO: Data Protection Officer)
\end{choices}

\question

Suponha que \( T \) seja uma árvore AVL inicialmente vazia, e considere a inserção dos elementos \(10, 20, 30, 5, 15, 2\) em \( T \), nesta ordem. Qual das sequências abaixo corresponde a um percurso de \( T \) em pré-ordem:
\begin{choices}
\choice 20, 10, 5, 2, 15, 30
\choice 10, 5, 2, 20, 15, 30
\choice 30, 20, 15, 10, 5, 2
\choice 15, 10, 5, 2, 20, 30
\choice 2, 5, 10, 15, 20, 30
\end{choices}

\question

Ao executar o seguinte comando SQL, nenhum dado foi retornado, ou seja, não

ocorreu nenhum erro de sintaxe mas nenhum dado foi encontrado:

\begin{tcolorbox}

 \begin{verbatim}1   SELECT c.nome AS cliente

 2   FROM clientes c

 3   LEFT JOIN pedidos p ON (p.cliente_id = c.cliente_id)

 4   WHERE p.cliente_id IS NULL

 5   ORDER BY cliente;

 \end{verbatim}

\end{tcolorbox}

Por que esse comando SQL não retornou nenhum dado?


\begin{choices}
\choice A junção que deveria ter sido realizada, na linha 3, era uma RIGHT JOIN ao invés de uma LEFT JOIN.
\choice A junção que deveria ter sido realizada, na linha 3, era uma FULL JOIN ao invés de uma RIGHT JOIN.
\choice Há um erro semântico na linha 3
\choice Todos os clientes cadastrados fizeram, pelo menos, um pedido de compra.
\choice Há um erro lógico na linha 4.
\end{choices}

\question

Assinale a alternativa INCORRETA:


\begin{choices}
\choice Os serviços orientados a conexões são sempre confiáveis garantindo a entrega ordenada e completa dos dados transmitidos.
\choice Nos serviços orientados a conexões há a necessidade de estabelecimento de uma conexão antes da transferência dos dados.
\choice O controle de fluxo tem como objetivo garantir que nenhum dos parceiros de uma comunicação inunda o outro enviando pacotes mais rápido do que ele pode tratar.
\choice Os serviços orientados a conexão podem ajudar no controle de congestionamento através da diminuição da taxa de transmissão durante um congestionamento em andamento.
\choice Serviços orientados a conexão podem ser implementados em subredes que funcionam no modo datagrama.
\end{choices}

\question

Uma característica de uma arquitetura RISC é:
\begin{choices}
\choice O processador é formado por válvulas e transistores
\choice Possui um pequeno conjunto de instruções simples
\choice Uma arquitetura de alto risco pois o mercado de hardware evolui muito rapidamente
\choice Possui um grande conjunto de instruções complexas
\choice A arquitetura é constituída de milhares de processadores
\end{choices}

\question

Um compilador detecta:
\begin{choices}
\choice erros nos resultados gerados pelo programa
\choice erros de sintaxe do programa 
\choice todos os erros citados acima
\choice erros aritméticos
\choice erros que podem ocorrer durante a execução do programa
\end{choices}

\question

Considere as seguintes afirmações sobre resolução de problemas em IA.

\begin{enumerate}[label=\Roman*.]

    \item A* é um conhecido algoritmo de busca heurística.

    \item O Minimax é um dos principais algoritmos para jogos de dois jogadores, como o xadrez.

    \item Busca em espaço de estados é uma das formas de resolução de problemas em IA.

\end{enumerate}

São corretas:
\begin{choices}
\choice Apenas I e III
\choice I, II e III
\choice Apenas III
\choice Apenas II e III
\choice Apenas I e II
\end{choices}

\question

Três atletas \( A \), \( B \) e \( C \) competiram, ao pares, numa corrida de \( d \) metros. Considerando que cada atleta teve o mesmo desempenho (ou seja, a mesma velocidade) ao competir com adversários distintos, e sabendo-se que



\begin{itemize}

    \item \( A \) venceu \( B \) chegando 20 metros à frente

    \item \( B \) venceu \( C \) chegando 10 metros à frente

    \item \( A \) venceu \( C \) chegando 28 metros à frente,

\end{itemize}



podemos afirmar que a corrida tem
\begin{choices}
\choice 100 metros
\choice 110 metros
\choice 50 metros
\choice 200 metros
\choice 150 metros
\end{choices}

\question

Assinale o argumento válido, onde \( S_1 \), \( S_2 \) indicam premissas e \( S \) a conclusão:


\begin{choices}
\choice \( S_1 \): Se o cavalo estiver cansado então ele perderá a corrida \\

              \( S_2 \): O cavalo ganhou a corrida \\

              \( S \): O cavalo estava descansado
\choice nenhuma das anteriores
\choice \( S_1 \): Se o cavalo estiver cansado então ele perderá a corrida \\

              \( S_2 \): O cavalo estava descansado \\

              \( S \): O cavalo ganhou a corrida
\choice \( S_1 \): Se o cavalo estiver cansado então ele perderá a corrida \\

              \( S_2 \): O cavalo estava descansado \\

              \( S \): O cavalo perdeu a corrida
\choice Se o cavalo estiver cansado então ele perderá a corrida \\

              \( S_2 \): O cavalo perdeu a corrida \\

              \( S \): O cavalo estava cansado
\end{choices}

\question

Considere o seguinte código para implementar exclusão mútua entre dois processos \(i\) e \(j\): 

Processo \(P_i\)

\begin{lstlisting}[language=C, basicstyle=\ttfamily\small]

do {

    while (turn != i); 

    // secao critica

    secao critica

    turn = j;

    // saída da secao critica

    codigo restante

} while (1);

\end{lstlisting} 
\begin{choices}
\choice A implementação garante exclusão mútua.
\choice Um processo bloqueia o outro mesmo não estando na seção crítica.
\choice Exige alternância estrita. 
\choice A implementação garante progresso. 
\choice Os processos fazem espera ativa.
\end{choices}

\question

Avalie o banco de dados de professores e alunos, mostrado a seguir. Esse banco

de dados pretende armazenar todas as turmas, sendo que um professor pode dar

aula para vários alunos ao mesmo tempo, e um mesmo aluno pode ter aula com

diferentes professores ao mesmo tempo.

\begin{figure}[H]

  \centering

  \caption{Professores, turmas e alunos}

  \label{fig:professores}

  %\fbox{

    \includegraphics[width=0.5\linewidth]{static/imgs/defaul.png}

  %}\\

  %\footnotesize{Fonte: xxxx.}

\end{figure}

Esse banco de dados resolveráo problema de armazenar os dados das turmas?
\begin{choices}
\choice Não, pois o relacionamento entre 'turmas' e 'professores' está errado: a PK 'cod\_turma' deveria ter ido para a tabela 'professores' como uma FK, mas no diagrama está ao contrário.
\choice Sim, pois o diagrama mostra claramente que um professor pode ter zero ou mais turmas, e cada aluno pode participar de zero ou mais turmas também.
\choice Sim, pois as cardinalidades e os sentidos dos relacionamentos entre 'professores' e 'alunos' indica, no final, um relacionamento N:N entre essas duas tabelas.
\choice Não, pois as cardinalidades e os sentidos dos relacionamentos entre 'professores' e 'turmas', e entre 'turmas' e 'alunos' não cria um relacionamento N:N entre 'professores' e 'alunos'.
\choice Não é possível determinar somente avaliando o diagrama exibido, seriam necessárias mais informações do dicionário de dados.
\end{choices}

\question

Em relação ao teste de software, qual das afirmações a seguir é INCORRETA: 
\begin{choices}
\choice O processo de depuração é a parte mais imprevisível do processo de teste pois um erro pode demorar uma hora, um dia ou um mês para ser diagnosticado e corrigido. 
\choice Os dados compilados quando a atividade de teste é levada a efeito proporcionam uma boa indicação da confiabilidade do software e alguma indicação da qualidade do software como um todo. 
\choice Um teste bem sucedido é aquele que revela um erro ainda não descoberto. 
\choice Um bom caso de teste é aquele que tem uma elevada probabilidade de revelar um erro ainda não descoberto. 
\choice A atividade de teste é o processo de executar um programa com a intenção de demonstrar a ausência de erros. 
\end{choices}

\question

Sobre a UML, quais das seguintes afirmações são verdadeiras?



\begin{itemize}

    \item[I)] A UML é o método de desenvolvimento de software mais utilizado na atualidade.

    \item[II)] A UML é uma evolução das linguagens para especificação dos conceitos dos métodos de Booch, OMT e OOSE e também de outros métodos de especificação de requisitos de software orientados a objetos ou não.

    \item[III)] A UML é composta dos seguintes diagramas: Diagrama de Caso de Uso, Diagrama de Classes, Diagrama de Colaboração, Diagrama de Estados, entre outros.

    \item[IV)] Em UML pode-se representar tão somente relacionamentos de Agregação, Associação e Composição.

\end{itemize}
\begin{choices}
\choice Nenhuma delas.
\choice Apenas as alternativas III e IV.
\choice Apenas as alternativas I, II e III.
\choice Apenas as alternativas II e III.
\choice Todas as alternativas.
\end{choices}

\question

A medida da interconexão entre os módulos de uma estrutura de software é denominada e que também é usada em projetos orientados a objetos é: 
\begin{choices}
\choice unidade funcional
\choice ocultamento da informação
\choice acoplamento
\choice abstração procedimental
\choice coesão
\end{choices}

\question

Na criptografia com chave pública:


\begin{choices}
\choice Para assinar digitalmente uma mensagem codifica-se a mesma com a chave pública do destinatário e esta é decifrada com a chave privada do destinatário.
\choice O sigilo é obtido através da codificação com a chave privada do destinatário e decifragem com a chave pública do destinatário.
\choice O sigilo é obtido através da codificação com a chave pública do destinatário e decifragem com a chave privada do destinatário.
\choice Para assinar digitalmente uma mensagem codifica-se a mesma com a chave pública do remetente e esta é decifrada com a chave privada do destinatário.
\choice O sigilo é obtido através da codificação com a chave privada do remetente e decifragem com a chave pública do destinatário.
\end{choices}

\question

O menor número possível de arestas em um grafo conexo com \( n \) vértices é:
\begin{choices}
\choice \( n - 1 \)
\choice \( n \)
\choice \( n/2 \)
\choice 1
\choice \( n^2 \)
\end{choices}

\question

O conjunto básico de atividades e a ordem em que são realizadas no processo de construção de um software definem o que é habitualmente denominado de ciclo de vida do software. O ciclo de vida tradicional (também denominado waterfall) ainda é hoje em dia um dos mais difundidos e tem por característica principal: 
\begin{choices}
\choice a avaliação constante dos resultados intermediários feita pelo cliente.
\choice o uso de formalização rigorosa em todas as etapas de desenvolvimento; 
\choice a priorização da análise dos riscos do desenvolvimento
\choice a abordagem sistemática para realização das atividades do desenvolvimento de software de modo que elas seguem um fluxo sequencial; 
\choice a codificação de uma versão executável do sistema desde as fases iniciais do desenvolvimento, de modo que o sistema final é incrementalmente construído, daí a alusão à ideia de "cascata" (waterfall); 
\end{choices}

\question

Considere as seguintes tabelas em uma base de dados relacional:\\

\textbf{Departamento (CodDepto, NomeDepto)}\\

\textbf{Empregado (CodEmp, NomeEmp, CodDepto)}\\

Deseja-se obter uma tabela na qual cada linha é a concatenação de uma linha da tabela Departamento com uma linha da tabela de Empregado. Caso um departamento não possua empregados, sua linha no resultado deve conter vazio (NULL) nos campos referentes ao empregado. A operação de álgebra relacional que deve ser aplicada para combinar estas duas tabelas é:
\begin{choices}
\choice Junção interna
\choice Divisão
\choice Projeção
\choice Junção externa
\choice União
\end{choices}

\question

A função abaixo computa a soma dos \( n \) primeiros números inteiros não negativos:



\begin{verbatim}

function sum(n: integer): integer;

begin

    if n = 0 then sum := 0

    else -------------

end;

\end{verbatim}



A parte que falta para completar a condição \texttt{else} é:
\begin{choices}
\choice \texttt{sum := (n-1) + sum(n)}
\choice \texttt{sum := n + sum(n-1)}
\choice \texttt{sum := n + sum(n)}
\choice \texttt{while n<>0 sum := sum + sum(n+1)}
\choice \texttt{sum := (n-1) + sum(n-1)}
\end{choices}

\question

Histograma de uma imagem com \( K \) tons de cinza é:
\begin{choices}
\choice Contagem dos pixels da imagem.
\choice Nenhuma alternativa acima.
\choice Contagem do número de tons de cinza que ocorreram na imagem.
\choice Contagem do número de vezes que cada um dos \( K \) tons de cinza ocorreu na imagem.
\choice Contagem do número de objetos encontrados na imagem.
\end{choices}

\question

(Adaptado de ENADE 2011) Considere o diagrama abaixo, que ilustra um banco de

dados que descreve os empregados que trabalham nos departamentos de uma empresa:

\begin{figure}[H]

  \centering

  \caption{Empregados e departamentos}

  \vspace{-0.2cm}

  \label{fig:empregados}

  %\fbox{

    \includegraphics[width=0.5\linewidth]{static/imgs/defaul.png}

  %}\\

  %\footnotesize{Fonte: xxxx.}

\end{figure}

\vspace{-0.4cm}

Na linguagem SQL, o comando mais simples para recuperar os códigos dos

departamentos que contém empregados que ganham menos do que 2000 e mais do que

5000 reais é:


\begin{choices}
\choice \begin{verbatim}SELECT cod_departamento

FROM empregados

WHERE salario <= 2000 AND salario >= 5000;

\end{verbatim}


\choice \begin{verbatim}SELECT cod_departamento

FROM empregados

WHERE salario BETWEEN 2000 AND 5000;

\end{verbatim}


\choice \begin{verbatim}SELECT cod_departametno

FROM empregados

WHERE salario <= 2000 OR salario >= 5000;

\end{verbatim}


\choice \begin{verbatim}SELECT cod_departamento

FROM empregados

WHERE salario < 2000 OR salario > 5000;

\end{verbatim}


\choice \begin{verbatim}SELECT cod_departamento

FROM empregados

WHERE salario < 2000 AND salario > 5000;

\end{verbatim}


\end{choices}

\question

Redes semânticas, frames e lógica são formalismos utilizados principalmente em: 
\begin{choices}
\choice descoberta de conhecimento em bases de dados
\choice IA distribuída
\choice inferência em sistemas especialistas
\choice representação de conhecimento
\choice redes neurais
\end{choices}

\question

Dentre os algoritmos de ordenação citados abaixo, qual é o que executa mais rápido para uma grande variedade de entrada de dados? 
\begin{choices}
\choice shellsort
\choice heapsort
\choice mergesort
\choice quicksort
\choice bolha
\end{choices}



    %%%%%%%%%%%%%%%%%%%%%%%%%%%%%%%%%%%%%%%%%%%%%%%%%%%%%%%%%%%%%%%%%%%%%%%%%%%%%%%%
    %%% Finaliza o bloco de questões:
    \end{questions}
    \end{document}
    